\documentclass[12pt,a4paper]{article}

% ── Paquetes ──────────────────────────────────────────────────────────────────
\usepackage[utf8]{inputenc}
\usepackage[T1]{fontenc}
\usepackage[spanish]{babel}
\usepackage{amsmath, amssymb}
\usepackage{booktabs}
\usepackage{array}
\usepackage[margin=2.5cm]{geometry}
\usepackage[most]{tcolorbox}

% ── Entornos de cajas ─────────────────────────────────────────────────────────
\tcbuselibrary{breakable}

\newtcolorbox{intuicion}{
  colback=blue!5!white,
  colframe=blue!50!black,
  fonttitle=\bfseries,
  title={Intuición},
  breakable
}

\newtcolorbox{interpretacion}{
  colback=green!5!white,
  colframe=green!50!black,
  fonttitle=\bfseries,
  title={Interpretación},
  breakable
}

% ── Metadatos ─────────────────────────────────────────────────────────────────
\title{\textbf{Guía de Práctica N$^\circ$5 — Política de Dividendos}\\[6pt]
       \large Ejercicio 1: Empresa REUNISA}
\author{Economía Financiera 1 \\ Universidad de Lima}
\date{Junio 2026}

% ══════════════════════════════════════════════════════════════════════════════
\begin{document}
\maketitle

\textbf{Datos del ejercicio:} En 2024, la empresa REUNISA distribuyó un
\textbf{dividendo trimestral (quarterly dividend)} de \$0.15 por acción.

% ──────────────────────────────────────────────────────────────────────────────
\subsection*{a) Emparejamiento de fechas}

\begin{intuicion}
Se pide identificar a qué etapa del proceso de distribución de dividendos
corresponde cada fecha. En finanzas, el pago de dividendos sigue una
secuencia cronológica fija: primero se anuncia, luego se define quién
tiene derecho, y finalmente se paga.
\end{intuicion}

La secuencia estándar es:
\[
\text{Anuncio} \;\longrightarrow\;
\text{Última fecha con dividendo} \;\longrightarrow\;
\text{Ex-Dividendo} \;\longrightarrow\;
\text{Registro/Cierre} \;\longrightarrow\;
\text{Pago}
\]

\begin{center}
\begin{tabular}{cll}
\toprule
\textbf{Fecha} & \textbf{Denominación} & \textbf{Rol en el proceso} \\
\midrule
31/10/24 & \textbf{Fecha del anuncio} (Declaration Date)
         & La junta declara el dividendo \\
24/11/24 & \textbf{Última fecha con dividendo} (Cum-Dividend Date)
         & Último día para comprar y cobrar \\
25/11/24 & \textbf{Fecha Ex-Dividendo} (Ex-Dividend Date)
         & Desde aquí el comprador no cobra \\
28/11/24 & \textbf{Fecha de cierre} (Record Date)
         & La empresa registra a los accionistas \\
12/12/24 & \textbf{Fecha de pago} (Payment Date)
         & Se acredita el dividendo \\
\bottomrule
\end{tabular}
\end{center}

\begin{interpretacion}
El emparejamiento correcto difiere de la disposición original del enunciado.
La clave lógica es que el anuncio (31/10) ocurre semanas antes del pago
(12/12), y la fecha ex-dividendo (25/11) sigue inmediatamente a la última
fecha con dividendo (24/11), de acuerdo con la liquidación T+2 vigente en la
mayoría de mercados.
\end{interpretacion}

% ──────────────────────────────────────────────────────────────────────────────
\subsection*{b) Fecha en que baja el precio}

\begin{intuicion}
Se pide identificar en qué fecha el mercado ajusta el precio de la acción
para reflejar que el derecho al dividendo ya no está incluido. Esto ocurre
el día en que la acción deja de ``llevar'' el dividendo consigo.
\end{intuicion}

El precio cae en la \textbf{fecha ex-dividendo (Ex-Dividend Date)},
correspondiente al \textbf{25/11/24}. A partir de ese día, quien compra
la acción ya no tiene derecho al dividendo declarado, por lo que el mercado
descuenta su valor:
\[
P_{\text{ex}} \;\approx\; P_{\text{cum}} - D
\]

\begin{interpretacion}
La caída esperada es de aproximadamente \$0.15 (el dividendo trimestral).
Este ajuste no representa una pérdida para el accionista que ya tenía la
acción: pierde precio pero gana el dividendo en efectivo. La riqueza total
no cambia en un mercado perfecto (Modigliani--Miller).
\end{interpretacion}

% ──────────────────────────────────────────────────────────────────────────────
\subsection*{c) Rentabilidad previsible por dividendo}

\begin{intuicion}
Se pide calcular el \textbf{rendimiento por dividendo (Dividend Yield)},
que mide cuánto representa el flujo anual de dividendos en relación al
precio de la acción. Es una métrica de retorno para el inversor que busca
ingresos periódicos.
\end{intuicion}

\textbf{Paso 1:} Calcular el dividendo anual.

Con 4 pagos trimestrales al año:
\[
D_{\text{anual}} = D_{\text{trimestral}} \times 4 = \$0.15 \times 4 = \$0.60
\]

\textbf{Paso 2:} Aplicar la fórmula del rendimiento por dividendo.
\[
\text{Dividend Yield} = \frac{D_{\text{anual}}}{P_0}
= \frac{\$0.60}{\$28.00}
= 0.02143
\]

\[
\textbf{Dividend Yield} \approx \textbf{2.14\%}
\]

\begin{interpretacion}
Por cada dólar invertido en REUNISA al precio de \$28, el inversor recibe
aproximadamente \$0.0214 anuales en dividendos. Este rendimiento modesto
sugiere que el mercado valora las expectativas de crecimiento de capital
más que el flujo de dividendos corriente.
\end{interpretacion}

% ──────────────────────────────────────────────────────────────────────────────
\subsection*{d) Ratio de distribución de dividendos}

\begin{intuicion}
El \textbf{ratio de distribución (Payout Ratio)} indica qué fracción de las
ganancias generadas por la empresa se devuelve a los accionistas como
dividendo. Un ratio alto implica generosidad con el accionista; uno bajo,
preferencia por reinvertir.
\end{intuicion}

\textbf{Datos:}
\begin{itemize}
    \item Dividendo anual por acción: $D_{\text{anual}} = \$0.60$
    \item \textbf{Utilidad por acción} (Earnings Per Share --- EPS)$= \$1.25$
\end{itemize}

\textbf{Fórmula:}
\[
\text{Payout Ratio} = \frac{D_{\text{anual}}}{\text{EPS}}
= \frac{\$0.60}{\$1.25}
= 0.48
\]

\[
\textbf{Payout Ratio} = \textbf{48\%}
\]

\begin{interpretacion}
REUNISA distribuye el 48\% de sus ganancias como dividendo, reteniendo el
52\% restante. Esta política moderada señala que la empresa equilibra la
retribución al accionista con la necesidad de autofinanciar su crecimiento.
Un Payout Ratio cercano al 50\% es típico de empresas maduras con flujos
de caja estables.
\end{interpretacion}

% ──────────────────────────────────────────────────────────────────────────────
\subsection*{e) Caída esperada por dividendo en acciones}

\begin{intuicion}
El \textbf{dividendo en acciones (Stock Dividend)} no es un pago en efectivo:
la empresa emite nuevas acciones y las distribuye proporcionalmente. Al haber
más acciones en circulación pero el mismo valor total de la empresa, el precio
por acción debe caer para que la riqueza total del accionista no cambie.
\end{intuicion}

\textbf{Datos:}
\begin{itemize}
    \item Precio actual: $P_0 = \$28$
    \item Tasa de \textbf{dividendo en acciones (Stock Dividend)}: $s = 10\% = 0.10$
\end{itemize}

\textbf{Paso 1:} Calcular el nuevo precio ajustado.

Tras el stock dividend, hay un 10\% más de acciones. El valor total de la
empresa no cambia, por lo que:
\[
P_{\text{nuevo}} = \frac{P_0}{1 + s} = \frac{\$28}{1.10} \approx \$25.45
\]

\textbf{Paso 2:} Calcular la caída esperada.
\[
\Delta P = P_0 - P_{\text{nuevo}} = \$28.00 - \$25.45 = \$2.55
\]

\[
\textbf{Caída esperada en el precio} \approx \textbf{\$2.55 por acción}
\]

\begin{interpretacion}
El precio cae de \$28 a \$25.45, una reducción del $9.09\%$. Sin embargo,
el accionista no pierde riqueza: si tenía 10 acciones a \$28 (\$280 en
total), ahora tiene 11 acciones a \$25.45 (\$280 en total). El stock
dividend es esencialmente un \textit{split} parcial: aumenta la liquidez
de la acción y puede tener efectos de señalización positivos, pero en un
mercado perfecto no crea ni destruye valor.
\end{interpretacion}

\end{document}
